\begin{lemma}[Boundary structure after a parallel extension]
Let \(n=q+1\), let \(T_n=\mathcal T.\mathrm{tree}(q+1)\), and let
\(\varphi_n=\mathcal T.\mathrm{coloring}(q+1)\).  Suppose the certified
current step is a PE, with exact connecting-edge set \(F_n=\{f_n\}\),
stored colors \(\delta_n,\gamma_n\), and supporting vertex \(v_n\).  Then
\[
 \partial_{\varphi_n,\gamma_n}(T_n)=\{f_n\}.
\]
For every identity \(g\in\partial_{\varphi_n,\delta_n}(T_n)\), there is a
tree vertex \(w\prec_{T_n}v_n\) incident with \(g\).  Finally, for every
\(\alpha\in\overline\varphi_n(T_n)\) with
\(\alpha\ne\delta_n\),
\[
 \partial_{\varphi_n,\alpha}(T_n)=\varnothing.
\]
All three assertions concern exact edge identities, so the singleton in
the first assertion is not merely a singleton endpoint pair.
\end{lemma}
\begin{proof}
Use the current certificate supplied by \noderef{TashkinovSeries}.  The
equality \(\Theta_n=\mathrm{PE}\) rules out the SE and RE constructors of
\noderef{TashkinovSeriesStep}; injectivity of \(\mathrm{some}\) and of
singleton finite sets identifies the constructor's witnesses with the
given \(f_n,\delta_n,\gamma_n,v_n\).  Write \(\pi\) for its
maximum-defective selected coloring and \(\rho\) for its second stable
coloring.

We first collect the pre-switch boundary information.  By
\noderef{MaximumDefectiveChoice}, \(\pi\) is
\((T_n,D_{n-1},\varphi_{n-1})\)-stable, \(v_n\) is maximum among all
defective vertices arising under such stable colorings, and \(f_n\) is a
defective boundary identity incident with \(v_n\), with
\(\pi(f_n)=\delta_n\).  Expanding \noderef{DefectiveEdge} identifies
\(\delta_n\) as the defective color of that boundary identity.  The PE
constructor says that \(\rho\) is stable
relative to the enlarged exact set \(D_{n-1}\cup\{\delta_n\}\).
Consequently every incident identity colored \(\delta_n\) under either
\(\pi\) or \(\rho\) has the same value under both colorings.  Formally,
one direction is the edge clause of \noderef{StableColoring}; the reverse
direction follows by symmetry in \noderef{sc2}, since both colorings are
proper and leave exactly \(e_0\) uncolored.  Thus their exact
\(\delta_n\)-boundary sets coincide, and in particular
\(\rho(f_n)=\delta_n\).

The same stability argument transfers closedness from
\(\varphi_{n-1}\) to \(\pi\) and then to \(\rho\).  Indeed, the PE
constructor states that \(V(T_n)\) is closed under \(\varphi_{n-1}\).
The vertex clause of \noderef{StableColoring} identifies the missing
colors at every tree vertex.  If an identity on the boundary acquired one
of those missing colors, symmetry from \noderef{sc2} would force its old
value to be the same, contradicting \noderef{ColoredVertexSetClosed}.
Hence every color in \(\overline\rho(T_n)\), including
\(\gamma_n\), has empty \(\rho\)-boundary.

Next consider a \(\delta_n\)-boundary identity \(g\) under \(\rho\), and
let \(w\in V(T_n)\) be its inside endpoint, supplied by
\noderef{BoundaryColorEdges} and \noderef{BoundaryEdges}; its relation to
\(g\) is the incidence
relation \noderef{Incident}.  The color \(\delta_n\) is defective under
\(\pi\), because \(f_n\) is a defective edge with that color; the equality
of the two \(\delta_n\)-boundary sets therefore makes \(w\) a defective
vertex under \(\pi\) by \noderef{DefectiveVertex}.  Maximality in
\noderef{MaximumDefectiveChoice} gives either \(w=v_n\) or
\(w\prec_{T_n}v_n\), where the latter relation is exactly
\noderef{TreeVertexOrder}.  The edgewise condition
\noderef{IsProperEdgeColoring} permits at most one
\(\delta_n\)-identity incident with \(v_n\), and that identity is
\(f_n\).  We have proved
\[
 g=f_n\quad\hbox{or}\quad
 \text{the inside endpoint of }g\text{ precedes }v_n.             \tag{1}
\]

We now identify the switched component.  The PE color rule
\noderef{ParallelExtensionColor} says that \(\gamma_n\) is missing at
\(v_n\) under \(\pi\), and stability gives the same under \(\rho\).
Since \(f_n\) is incident there with color \(\delta_n\), properness makes
\(f_n\) the unique two-color edge at \(v_n\); in particular
\(\gamma_n\ne\delta_n\).  The \noderef{TreeExit} witness supplies a
simple path from an outside vertex to \(v_n\), meeting \(T_n\) only in
\(v_n\).  Its outside endpoint misses one of the two colors by
\noderef{TreeExitPath} and has an incident path edge of the other color,
so both endpoints of the path have incidence-graph degree one.  All other
nodes have degree at most two by
\noderef{TwoColorIncidenceDegreeAtMostTwo}.  If the connected component
had a node outside this path, a shortest route from it to the path would
give either a third neighbor at an internal path node or a second neighbor
at an endpoint.  This is impossible.  Thus the path is the whole
\noderef{TwoColorComponent}; it meets the tree only at \(v_n\), and
\(f_n\) is its unique boundary identity.

The definition \noderef{TwoColorComponentSwap} now gives the boundary
calculation edge by edge.  Before the switch, the component's unique
boundary identity \(f_n\) has color \(\delta_n\), and the closedness just
proved says that there is no \(\gamma_n\)-boundary identity.  The switch
changes \(f_n\) to \(\gamma_n\), changes no other boundary identity to
\(\gamma_n\), and removes \(f_n\) from the \(\delta_n\)-boundary.
Therefore
\[
 \partial_{\varphi_n,\gamma_n}(T_n)=\{f_n\}.                     \tag{2}
\]
Every remaining \(\delta_n\)-boundary identity is a
\(\delta_n\)-boundary identity of \(\rho\) different from \(f_n\), so
(1) gives an incident inside vertex strictly preceding \(v_n\).  This is
the second assertion.

For the last assertion, compare missing colors and boundary identities
before and after the same switch.  The full switched component meets
\(T_n\) only at its supporting endpoint.  Thus the only missing-color
replacement on the tree is the PE replacement of \(\gamma_n\) by
\(\delta_n\) at that endpoint, and the only boundary replacement is the
change of the unique identity \(f_n\) from \(\delta_n\) to
\(\gamma_n\).  In particular, a color
\(\alpha\in\overline\varphi_n(T_n)\) with
\(\alpha\ne\delta_n\) is a pre-switch tree-missing color whose boundary
identities are unchanged by the switch.  The pre-switch closedness proved
above makes that boundary empty.  Hence
\(\partial_{\varphi_n,\alpha}(T_n)=\varnothing\), as required.
\end{proof}
